\section{The Configuration File}
\label{sec:parameters}

\subsection{How the file works}

A study is described by one YAML file (\texttt{configs/weisshorn.yaml} is the complete, commented example). YAML
is a plain-text format: \emph{sections} contain \emph{parameters} written \texttt{name: value}, indented by two
spaces; lists are written \texttt{[a, b, c]}; everything after \texttt{\#} is a comment.
\begin{lstlisting}[language=yaml]
vehicle:
  diameter: 0.180        # m
  mass_wet: 75.0         # kg
nose:
  profiles: [ogive, power050, vonkarman]
  fineness: [4.0, 4.5, 5.0]
\end{lstlisting}

\begin{itemize}[leftmargin=*,itemsep=1pt]
\item \textbf{Only what differs from the default is needed.} Every parameter has a default
(Section~\ref{sec:paramtables}); a file with only a \texttt{vehicle} section is valid.
\item \textbf{Units are SI} (m, kg, s, N, K) unless the table says otherwise (\textdegree, \textdegree C).
Write lengths in metres: \texttt{diameter: 0.180}, not \texttt{180}.
\item \textbf{Relative paths} (database, output folder, CFD folders, plugins) are relative to the file that contains them.
\item \textbf{Inheritance}: a file that starts with \texttt{base: weisshorn.yaml} takes every value of that file
and overrides only what it lists (both worked examples of Section~\ref{sec:running} do this). Sections are merged
parameter by parameter; a list given in the new file replaces the inherited list.
\item \textbf{Checking}: \texttt{python -m nosecone validate my.yaml} reads the file and stops at the first problem
with a precise message, e.g.\ \texttt{unknown parameter 'nose.profils' (did you mean 'profiles'?)} or
\texttt{'launch.rail\_elevation' = 95.0 is outside the valid range (0, 90]}. Add \texttt{-{}-show} to print every
value the run will use (defaults included). The same file is saved next to the results as
\texttt{config\_used.yaml}.
\item \textbf{Plugins}: \key{project.plugins} lists Python files imported before the run; they can add profile
families (Section~\ref{sec:newfamily}).
\end{itemize}

\subsection{All parameters}
\label{sec:paramtables}

The tables below are generated from the code (\texttt{python -m nosecone describe-config -{}-latex}), so they always
match the installed version. ``Valid'' gives the accepted range or values; parameters marked $^\ast$ are advanced
numerical settings that normally stay at their default. Parameters whose default is ``see text'' are structured
values, described in Section~\ref{sec:structured}.

\subsection*{\texttt{project}: Project}
\addcontentsline{toc}{subsection}{\texttt{project}: Project}
{\footnotesize
\begin{longtable}{@{}>{\raggedright\arraybackslash}p{3.3cm}>{\raggedright\arraybackslash}p{2.6cm}>{\raggedright\arraybackslash}p{1.3cm}>{\raggedright\arraybackslash}p{1.9cm}>{\raggedright\arraybackslash}p{6.1cm}@{}}
\toprule Parameter & Default & Unit & Valid & Meaning \\ \midrule \endhead \bottomrule \endfoot
\texttt{name} & \texttt{Weisshorn} & -- & -- & Vehicle or study name, used in plot titles and file names \\
\texttt{output\_dir} & \texttt{output} & path & -- & Folder for all outputs (created if missing) \\
\texttt{plugins} & [\,] & path & -- & Python files imported before the run, e.g. new profile families (see the user guide) \\
\end{longtable}\addtocounter{table}{-1}}

\subsection*{\texttt{vehicle}: Vehicle}
\addcontentsline{toc}{subsection}{\texttt{vehicle}: Vehicle}
{\footnotesize
\begin{longtable}{@{}>{\raggedright\arraybackslash}p{3.3cm}>{\raggedright\arraybackslash}p{2.6cm}>{\raggedright\arraybackslash}p{1.3cm}>{\raggedright\arraybackslash}p{1.9cm}>{\raggedright\arraybackslash}p{6.1cm}@{}}
\toprule Parameter & Default & Unit & Valid & Meaning \\ \midrule \endhead \bottomrule \endfoot
\texttt{diameter} & 0.18 & m & > 0 & External (body) diameter D; the nose base diameter \\
\texttt{length} & 5.4 & m & > 0 & Total vehicle length including the reference nose \\
\texttt{reference\_nose\_length} & 0.72 & m & > 0 & Nose length included in `length`; the body is length - reference\_nose\_length \\
\texttt{mass\_wet} & 75 & kg & > 0 & Lift-off mass \\
\texttt{mass\_dry} & 45 & kg & > 0 & Burn-out mass used by mass\_model 'lvo' \\
\texttt{surface\_roughness} & 2e-05 & m & >= 0 & Equivalent sand roughness of the skin (painted CFRP \~{} 20 um) \\
\texttt{nozzle\_exit\_diameter} & 0.075 & m & >= 0 & Nozzle exit diameter; removed from the base area while burning \\
\end{longtable}\addtocounter{table}{-1}}

\subsection*{\texttt{engine}: Engine}
\addcontentsline{toc}{subsection}{\texttt{engine}: Engine}
{\footnotesize
\begin{longtable}{@{}>{\raggedright\arraybackslash}p{3.3cm}>{\raggedright\arraybackslash}p{2.6cm}>{\raggedright\arraybackslash}p{1.3cm}>{\raggedright\arraybackslash}p{1.9cm}>{\raggedright\arraybackslash}p{6.1cm}@{}}
\toprule Parameter & Default & Unit & Valid & Meaning \\ \midrule \endhead \bottomrule \endfoot
\texttt{thrust} & 4500 & N & > 0 & Thrust (constant over the burn) \\
\texttt{burn\_time} & 9.1 & s & > 0 & Burn duration \\
\texttt{total\_impulse} & 40960 & N s & > 0 & Total impulse (sets the propellant mass for mass\_model 'isp') \\
\texttt{specific\_impulse} & 200 & s & > 0 & Specific impulse (propellant mass = I\_t / (I\_sp g0)) \\
\end{longtable}\addtocounter{table}{-1}}

\subsection*{\texttt{fins}: Fins}
\addcontentsline{toc}{subsection}{\texttt{fins}: Fins}
{\footnotesize
\begin{longtable}{@{}>{\raggedright\arraybackslash}p{3.3cm}>{\raggedright\arraybackslash}p{2.6cm}>{\raggedright\arraybackslash}p{1.3cm}>{\raggedright\arraybackslash}p{1.9cm}>{\raggedright\arraybackslash}p{6.1cm}@{}}
\toprule Parameter & Default & Unit & Valid & Meaning \\ \midrule \endhead \bottomrule \endfoot
\texttt{count} & 4 & -- & [1, 8] & Number of fins \\
\texttt{root\_chord} & 0.3 & m & > 0 & Root chord \\
\texttt{tip\_chord} & 0.12 & m & >= 0 & Tip chord \\
\texttt{span} & 0.16 & m & > 0 & Semi-span (root to tip) \\
\texttt{sweep} & 0.18 & m & >= 0 & Leading-edge sweep length (root LE to tip LE, axial) \\
\texttt{thickness} & 0.005 & m & > 0 & Fin thickness \\
\texttt{edge\_fraction} & 0.15 & -- & (0, 0.5] & Length of each LE/TE wedge of the hexagonal section, as a fraction of the mean chord \\
\end{longtable}\addtocounter{table}{-1}}

\subsection*{\texttt{launch}: Launch}
\addcontentsline{toc}{subsection}{\texttt{launch}: Launch}
{\footnotesize
\begin{longtable}{@{}>{\raggedright\arraybackslash}p{3.3cm}>{\raggedright\arraybackslash}p{2.6cm}>{\raggedright\arraybackslash}p{1.3cm}>{\raggedright\arraybackslash}p{1.9cm}>{\raggedright\arraybackslash}p{6.1cm}@{}}
\toprule Parameter & Default & Unit & Valid & Meaning \\ \midrule \endhead \bottomrule \endfoot
\texttt{rail\_length} & 12 & m & > 0 & Launch rail length \\
\texttt{rail\_elevation} & 84 & \textdegree & (0, 90] & Rail elevation above the horizontal \\
\texttt{site\_altitude} & 0 & m & [-500, 5000] & Launch site altitude above sea level \\
\end{longtable}\addtocounter{table}{-1}}

\subsection*{\texttt{atmosphere}: Atmosphere}
\addcontentsline{toc}{subsection}{\texttt{atmosphere}: Atmosphere}
{\footnotesize
\begin{longtable}{@{}>{\raggedright\arraybackslash}p{3.3cm}>{\raggedright\arraybackslash}p{2.6cm}>{\raggedright\arraybackslash}p{1.3cm}>{\raggedright\arraybackslash}p{1.9cm}>{\raggedright\arraybackslash}p{6.1cm}@{}}
\toprule Parameter & Default & Unit & Valid & Meaning \\ \midrule \endhead \bottomrule \endfoot
\texttt{model} & \texttt{isa\_troposphere} & -- & isa\_troposphere & Atmosphere model \\
\texttt{ground\_temperature} & 288.15 & K & [200, 340] & Temperature at the launch site (288.15 K = ISA) \\
\texttt{ground\_pressure} & 101325 & Pa & > 0 & Pressure at the reference altitude of the model \\
\texttt{lapse\_rate} & 0.0065 & K/m & > 0 & Temperature lapse rate in the troposphere \\
\texttt{tropopause\_temperature} & 216.65 & K & > 0 & Temperature floor (tropopause / lower stratosphere) \\
\end{longtable}\addtocounter{table}{-1}}

\subsection*{\texttt{flight}: Flight model}
\addcontentsline{toc}{subsection}{\texttt{flight}: Flight model}
{\footnotesize
\begin{longtable}{@{}>{\raggedright\arraybackslash}p{3.3cm}>{\raggedright\arraybackslash}p{2.6cm}>{\raggedright\arraybackslash}p{1.3cm}>{\raggedright\arraybackslash}p{1.9cm}>{\raggedright\arraybackslash}p{6.1cm}@{}}
\toprule Parameter & Default & Unit & Valid & Meaning \\ \midrule \endhead \bottomrule \endfoot
\texttt{mass\_model} & \texttt{isp} & -- & isp, lvo & Nominal mass model: 'isp' (propellant = I\_t/(I\_sp g0)) or 'lvo' (wet -> dry) \\
\texttt{calibrate} & \texttt{true} & -- & -- & Scale the non-nose drag by a factor k so that the baseline reaches target\_apogee \\
\texttt{target\_apogee} & 9000 & m & > 0 & Apogee the baseline nose must reach when calibrate is true \\
\texttt{drag\_factor} & 1 & -- & > 0 & Factor k on the non-nose drag when calibrate is false \\
\texttt{calibration\_bracket}\,$^\ast$ & \texttt{[0.3, 1.5]} & -- & -- & Search interval for k \\
\texttt{calibration\_tolerance}\,$^\ast$ & 0.0001 & -- & > 0 & Absolute tolerance on k \\
\texttt{t\_max}\,$^\ast$ & 120 & s & > 0 & Integration time limit (apogee normally ends the integration earlier) \\
\texttt{max\_step}\,$^\ast$ & 0.02 & s & > 0 & Maximum ODE time step \\
\texttt{rtol}\,$^\ast$ & 1e-08 & -- & > 0 & ODE relative tolerance \\
\texttt{atol}\,$^\ast$ & 1e-08 & -- & > 0 & ODE absolute tolerance \\
\texttt{output\_step}\,$^\ast$ & 0.05 & s & > 0 & Output sampling of the trade-study trajectories (sets M\_max) \\
\texttt{detail\_output\_step}\,$^\ast$ & 0.02 & s & > 0 & Output sampling of the trajectories used for heating and plots \\
\texttt{envelope\_mass\_model} & \texttt{lvo} & -- & isp, lvo & Mass model of the hot-day envelope trajectory \\
\texttt{envelope\_ground\_temperature} & 303.15 & K & [200, 340] & Ground temperature of the hot-day envelope \\
\texttt{mass\_sensitivity\_step}\,$^\ast$ & 1 & kg & > 0 & Mass step for the apogee sensitivity dz/dm \\
\texttt{cd\_sensitivity\_step}\,$^\ast$ & 0.01 & -- & > 0 & C\_D step for the apogee sensitivity dz/dC\_D \\
\end{longtable}\addtocounter{table}{-1}}

\subsection*{\texttt{nose}: Nose: design space and selected design}
\addcontentsline{toc}{subsection}{\texttt{nose}: Nose: design space and selected design}
{\footnotesize
\begin{longtable}{@{}>{\raggedright\arraybackslash}p{3.3cm}>{\raggedright\arraybackslash}p{2.6cm}>{\raggedright\arraybackslash}p{1.3cm}>{\raggedright\arraybackslash}p{1.9cm}>{\raggedright\arraybackslash}p{6.1cm}@{}}
\toprule Parameter & Default & Unit & Valid & Meaning \\ \midrule \endhead \bottomrule \endfoot
\texttt{profiles} & see text & -- & -- & Profiles compared in the sweep (built-in names or keys of custom\_profiles) \\
\texttt{fineness} & \texttt{[3, 3.5, 4, 4.5, 5]} & -- & -- & Fineness ratios L/D compared in the sweep \\
\texttt{baseline\_profile} & \texttt{vonkarman} & -- & -- & Reference nose for apogee and static-margin changes \\
\texttt{baseline\_fineness} & 4 & -- & > 0 & Fineness of the reference nose \\
\texttt{baseline\_label} & \texttt{baseline} & -- & -- & How the reference length is called in table captions (e.g. 'LVS length') \\
\texttt{design\_profile} & \texttt{vonkarman} & -- & -- & Selected profile (CAD export, heating, final drawing) \\
\texttt{design\_fineness} & 4.5 & -- & > 0 & Selected fineness L/D \\
\texttt{design\_tip\_radius\_ratio} & 0.05 & -- & [0, 0.5] & Selected spherical tip radius r\_n / R (0 = sharp) \\
\texttt{tip\_radius\_ratios} & \texttt{[0]} & -- & -- & Tip radii r\_n / R compared in the sweep (needs CFD data for each) \\
\texttt{shoulder\_length} & 0.06 & m & >= 0 & Cylindrical shoulder behind the curved part (coupler bond) \\
\texttt{custom\_profiles} & \{\,\} & -- & -- & Extra profiles: name -> \{family: ..., <parameter>: value, label: ...\} \\
\end{longtable}\addtocounter{table}{-1}}

\subsection*{\texttt{aerodynamics}: Aerodynamics (wave-drag model)}
\addcontentsline{toc}{subsection}{\texttt{aerodynamics}: Aerodynamics (wave-drag model)}
{\footnotesize
\begin{longtable}{@{}>{\raggedright\arraybackslash}p{3.3cm}>{\raggedright\arraybackslash}p{2.6cm}>{\raggedright\arraybackslash}p{1.3cm}>{\raggedright\arraybackslash}p{1.9cm}>{\raggedright\arraybackslash}p{6.1cm}@{}}
\toprule Parameter & Default & Unit & Valid & Meaning \\ \midrule \endhead \bottomrule \endfoot
\texttt{database} & \texttt{data/euler\_results.csv} & path & -- & CSV of Euler wave-drag results (non-dimensional, reusable for any diameter) \\
\texttt{user\_database} & (empty) & path & -- & Optional second CSV for cases computed with compute\_missing (empty: they are added to `database`) \\
\texttt{studies} & \texttt{[main, f45]} & -- & -- & Database studies used for the wave-drag model \\
\texttt{grid} & \texttt{medium} & -- & coarse, medium, fine & Grid level of the database rows used \\
\texttt{min\_fineness}\,$^\ast$ & 0 & -- & >= 0 & Database rows below this fineness are ignored \\
\texttt{max\_fineness}\,$^\ast$ & 5 & -- & > 0 & Database rows above this fineness are ignored \\
\texttt{machs} & \texttt{[1.1, 1.2, 1.4, 1.6]} & -- & -- & Supersonic Mach numbers at which the wave drag is tabulated \\
\texttt{transonic\_onset} & 0.8 & -- & [0.5, 1] & Mach number where the nose wave drag starts to rise \\
\texttt{transonic\_full} & 1 & -- & [0.8, 1.2] & Mach number where the rise reaches the value at transonic\_match \\
\texttt{transonic\_match} & 1.1 & -- & [1, 1.5] & Lowest Euler Mach number; the model is flat between transonic\_full and this \\
\texttt{compute\_missing} & \texttt{false} & -- & -- & Run the Euler solver for profiles / fineness / tip radii missing from the database \\
\texttt{missing\_machs} & \texttt{[1.6, 1.4, 1.2, 1.1]} & -- & -- & Mach sequence for cases computed on demand (highest first) \\
\end{longtable}\addtocounter{table}{-1}}

\subsection*{\texttt{payload}: Payload}
\addcontentsline{toc}{subsection}{\texttt{payload}: Payload}
{\footnotesize
\begin{longtable}{@{}>{\raggedright\arraybackslash}p{3.3cm}>{\raggedright\arraybackslash}p{2.6cm}>{\raggedright\arraybackslash}p{1.3cm}>{\raggedright\arraybackslash}p{1.9cm}>{\raggedright\arraybackslash}p{6.1cm}@{}}
\toprule Parameter & Default & Unit & Valid & Meaning \\ \midrule \endhead \bottomrule \endfoot
\texttt{enabled} & \texttt{true} & -- & -- & Evaluate the payload fit and nose-bay length \\
\texttt{name} & \texttt{3U} & -- & -- & Payload name used in labels \\
\texttt{length} & 0.3405 & m & > 0 & Payload length \\
\texttt{width} & 0.1 & m & > 0 & Payload width (square section), used in drawings \\
\texttt{required\_radius} & 0.077 & m & > 0 & Outer nose radius needed to house the payload section (half-diagonal + clearance + guide + shell) \\
\texttt{axial\_clearance} & 0.01 & m & >= 0 & Axial margin for the payload bulkhead \\
\end{longtable}\addtocounter{table}{-1}}

\subsection*{\texttt{structure}: Structure (mass model)}
\addcontentsline{toc}{subsection}{\texttt{structure}: Structure (mass model)}
{\footnotesize
\begin{longtable}{@{}>{\raggedright\arraybackslash}p{3.3cm}>{\raggedright\arraybackslash}p{2.6cm}>{\raggedright\arraybackslash}p{1.3cm}>{\raggedright\arraybackslash}p{1.9cm}>{\raggedright\arraybackslash}p{6.1cm}@{}}
\toprule Parameter & Default & Unit & Valid & Meaning \\ \midrule \endhead \bottomrule \endfoot
\texttt{shell\_areal\_density} & 0.98 & kg/m\textsuperscript{2} & > 0 & Shell mass per wetted area (Firehorn as-built: 0.98) \\
\texttt{tip\_mass} & 0.06 & kg & >= 0 & Machined tip insert mass (same for all profiles) \\
\texttt{shell\_cg\_fraction}\,$^\ast$ & 0.35 & -- & [0, 1] & Shell centre of gravity ahead of the shoulder, as a fraction of L \\
\end{longtable}\addtocounter{table}{-1}}

\subsection*{\texttt{heating}: Heating}
\addcontentsline{toc}{subsection}{\texttt{heating}: Heating}
{\footnotesize
\begin{longtable}{@{}>{\raggedright\arraybackslash}p{3.3cm}>{\raggedright\arraybackslash}p{2.6cm}>{\raggedright\arraybackslash}p{1.3cm}>{\raggedright\arraybackslash}p{1.9cm}>{\raggedright\arraybackslash}p{6.1cm}@{}}
\toprule Parameter & Default & Unit & Valid & Meaning \\ \midrule \endhead \bottomrule \endfoot
\texttt{enabled} & \texttt{true} & -- & -- & Run the heating analysis \\
\texttt{trajectory\_from} & \texttt{design} & -- & design, baseline & Nose used to compute the heating trajectories \\
\texttt{materials} & see text & -- & -- & Material properties: name -> \{label, short (optional), conductivity [W/(m K)], density [kg/m\^{}3], specific\_heat [J/(kg K)]\} \\
\texttt{shell\_material} & \texttt{CFRP} & -- & -- & Shell material (key of materials) \\
\texttt{shell\_thickness} & 0.0012 & m & > 0 & Shell thickness \\
\texttt{shell\_stations} & \texttt{[0.05 m, 0.5 L, 1.0 L]} & -- & -- & Shell stations: 'x m' from the tip or 'f L' as a fraction of L \\
\texttt{tip\_material} & \texttt{Al6082} & -- & -- & Tip insert material (lumped model) \\
\texttt{tip\_length} & 0.08 & m & > 0 & Axial length of the tip insert from the apex \\
\texttt{printed\_tip\_material} & \texttt{PETCF} & -- & -- & Material of the comparison printed tip (1D wall at the stagnation point) \\
\texttt{printed\_tip\_thickness} & 0.004 & m & > 0 & Wall thickness of the printed tip \\
\texttt{wall\_nodes}\,$^\ast$ & 40 & -- & [5, 400] & Nodes through the wall thickness \\
\texttt{bond\_limit} & 120 & \textdegree C & -- & Adhesive temperature limit drawn on the shell plot \\
\texttt{bond\_label} & \texttt{DP490 bond limit} & -- & -- & Label of the adhesive limit \\
\texttt{printed\_tip\_hdt} & \texttt{[70, 80]} & \textdegree C & -- & Heat-deflection temperature range of the printed tip polymer \\
\end{longtable}\addtocounter{table}{-1}}

\subsection*{\texttt{bluntness}: Tip bluntness}
\addcontentsline{toc}{subsection}{\texttt{bluntness}: Tip bluntness}
{\footnotesize
\begin{longtable}{@{}>{\raggedright\arraybackslash}p{3.3cm}>{\raggedright\arraybackslash}p{2.6cm}>{\raggedright\arraybackslash}p{1.3cm}>{\raggedright\arraybackslash}p{1.9cm}>{\raggedright\arraybackslash}p{6.1cm}@{}}
\toprule Parameter & Default & Unit & Valid & Meaning \\ \midrule \endhead \bottomrule \endfoot
\texttt{enabled} & \texttt{true} & -- & -- & Evaluate the tip-bluntness study \\
\texttt{study} & \texttt{blunt} & -- & -- & Database study with blunted noses \\
\texttt{profile} & \texttt{vonkarman} & -- & -- & Profile of the bluntness study \\
\texttt{fineness} & 4 & -- & > 0 & Fineness of the bluntness study \\
\texttt{report\_ratios} & \texttt{[0.05, 0.1]} & -- & -- & Tip radii r\_n/R whose drag and apogee cost are reported \\
\texttt{report\_mach} & 1.4 & -- & [1, 5] & Mach number at which the drag cost of the tip is reported \\
\texttt{radius\_range} & \texttt{[0.002, 0.02]} & m & -- & Tip radius range of the heat-transfer curve \\
\end{longtable}\addtocounter{table}{-1}}

\subsection*{\texttt{robustness}: Robustness}
\addcontentsline{toc}{subsection}{\texttt{robustness}: Robustness}
{\footnotesize
\begin{longtable}{@{}>{\raggedright\arraybackslash}p{3.3cm}>{\raggedright\arraybackslash}p{2.6cm}>{\raggedright\arraybackslash}p{1.3cm}>{\raggedright\arraybackslash}p{1.9cm}>{\raggedright\arraybackslash}p{6.1cm}@{}}
\toprule Parameter & Default & Unit & Valid & Meaning \\ \midrule \endhead \bottomrule \endfoot
\texttt{enabled} & \texttt{true} & -- & -- & Run the robustness study \\
\texttt{fineness} & 4.5 & -- & > 0 & Fineness at which all profiles are compared \\
\texttt{reference\_profile} & \texttt{vonkarman} & -- & -- & Apogees are shown relative to this profile \\
\texttt{band} & 30 & m & >= 0 & Half-width of the 'equivalent' band in the plot \\
\texttt{scenarios} & see text & -- & -- & Scenarios: name + any of transonic\_onset, transonic\_factor, drag\_factor\_scale, mass\_model, ground\_temperature \\
\end{longtable}\addtocounter{table}{-1}}

\subsection*{\texttt{trade\_matrix}: Decision matrix}
\addcontentsline{toc}{subsection}{\texttt{trade\_matrix}: Decision matrix}
{\footnotesize
\begin{longtable}{@{}>{\raggedright\arraybackslash}p{3.3cm}>{\raggedright\arraybackslash}p{2.6cm}>{\raggedright\arraybackslash}p{1.3cm}>{\raggedright\arraybackslash}p{1.9cm}>{\raggedright\arraybackslash}p{6.1cm}@{}}
\toprule Parameter & Default & Unit & Valid & Meaning \\ \midrule \endhead \bottomrule \endfoot
\texttt{enabled} & \texttt{true} & -- & -- & Build the weighted trade matrix \\
\texttt{fineness} & 4.5 & -- & > 0 & Fineness at which the matrix is built \\
\texttt{weights} & see text & -- & -- & Criterion weights (should add up to 1) \\
\texttt{manufacture} & see text & -- & -- & Judgement score per profile (1-5) \\
\texttt{heritage} & see text & -- & -- & Judgement score per profile (1-5) \\
\end{longtable}\addtocounter{table}{-1}}

\subsection*{\texttt{validation}: CFD validation plots}
\addcontentsline{toc}{subsection}{\texttt{validation}: CFD validation plots}
{\footnotesize
\begin{longtable}{@{}>{\raggedright\arraybackslash}p{3.3cm}>{\raggedright\arraybackslash}p{2.6cm}>{\raggedright\arraybackslash}p{1.3cm}>{\raggedright\arraybackslash}p{1.9cm}>{\raggedright\arraybackslash}p{6.1cm}@{}}
\toprule Parameter & Default & Unit & Valid & Meaning \\ \midrule \endhead \bottomrule \endfoot
\texttt{cone\_study} & \texttt{cones} & -- & -- & Database study of sharp cones (compared with Taylor-Maccoll) \\
\texttt{cone\_main\_fineness} & 4 & -- & > 0 & Cone of the main study added to the validation plot \\
\texttt{grid\_study} & \texttt{grid} & -- & -- & Database study of the grid-convergence runs \\
\texttt{grid\_profile} & \texttt{vonkarman} & -- & -- & Profile of the grid study \\
\texttt{grid\_fineness} & 4 & -- & > 0 & Fineness of the grid study \\
\texttt{grid\_machs} & \texttt{[1.4, 1.2]} & -- & -- & Mach numbers of the grid study \\
\end{longtable}\addtocounter{table}{-1}}

\subsection*{\texttt{cfd}: CFD solver and studies}
\addcontentsline{toc}{subsection}{\texttt{cfd}: CFD solver and studies}
{\footnotesize
\begin{longtable}{@{}>{\raggedright\arraybackslash}p{3.3cm}>{\raggedright\arraybackslash}p{2.6cm}>{\raggedright\arraybackslash}p{1.3cm}>{\raggedright\arraybackslash}p{1.9cm}>{\raggedright\arraybackslash}p{6.1cm}@{}}
\toprule Parameter & Default & Unit & Valid & Meaning \\ \midrule \endhead \bottomrule \endfoot
\texttt{cp\_dir} & \texttt{data/cfd/cp} & path & -- & Folder of the wall-pressure files (one .npz per case) \\
\texttt{fields\_dir} & \texttt{data/cfd/fields} & path & -- & Folder of the saved flow fields \\
\texttt{log} & \texttt{data/euler\_sweep\_log.csv} & path & -- & Run log appended by the solver \\
\texttt{workers} & 4 & -- & [1, 64] & Parallel processes \\
\texttt{cfl} & 0.85 & -- & [0.05, 1] & CFL number of the local time stepping \\
\texttt{fallback\_cfl}\,$^\ast$ & 0.5 & -- & [0.05, 1] & CFL of the restart from free stream if a run diverges \\
\texttt{max\_iter} & 6000 & -- & >= 100 & Iteration limit per Mach number \\
\texttt{fine\_iter\_factor}\,$^\ast$ & 2 & -- & [1, 10] & Iteration-limit multiplier on the fine grid \\
\texttt{tolerance} & 2e-05 & -- & > 0 & Convergence: relative change of C\_D,w over `check` iterations \\
\texttt{check}\,$^\ast$ & 250 & -- & >= 10 & Iterations between convergence checks \\
\texttt{first\_order\_iters}\,$^\ast$ & 300 & -- & >= 0 & First-order iterations at the start of the first Mach number \\
\texttt{grids} & see text & -- & -- & Grid levels: name -> [cells along nose, cells along cylinder, cells normal, first-cell height / R] \\
\texttt{cylinder\_length}\,$^\ast$ & 5 & R & > 0 & Cylinder behind the nose in the CFD domain \\
\texttt{upstream\_distance}\,$^\ast$ & 1 & R & > 0 & Distance of the outer boundary ahead of the tip (scaled by 1 + 0.3 L/D) \\
\texttt{outer\_angle}\,$^\ast$ & 72 & \textdegree & [30, 89] & Asymptotic angle of the outer boundary (hyperbola) \\
\texttt{save\_fields}\,$^\ast$ & \texttt{[1.4, 1.2]} & -- & -- & Mach numbers at which flow fields are saved (studies with save\_fields: true) \\
\texttt{studies} & [\,] & -- & -- & Euler studies: list of \{name, profiles, fineness, tip\_radius\_ratios, grid, machs, save\_fields\} \\
\end{longtable}\addtocounter{table}{-1}}

\subsection*{\texttt{cad}: CAD export}
\addcontentsline{toc}{subsection}{\texttt{cad}: CAD export}
{\footnotesize
\begin{longtable}{@{}>{\raggedright\arraybackslash}p{3.3cm}>{\raggedright\arraybackslash}p{2.6cm}>{\raggedright\arraybackslash}p{1.3cm}>{\raggedright\arraybackslash}p{1.9cm}>{\raggedright\arraybackslash}p{6.1cm}@{}}
\toprule Parameter & Default & Unit & Valid & Meaning \\ \midrule \endhead \bottomrule \endfoot
\texttt{profile\_csv} & \texttt{true} & -- & -- & Write the profile coordinates of the design (x from the tip, r) in mm \\
\texttt{profile\_points} & 200 & -- & [10, 100000] & Points of the profile CSV \\
\texttt{profile\_clustering}\,$^\ast$ & 1.6 & -- & [1, 4] & Clustering exponent towards the tip of the CSV points (1 = uniform) \\
\texttt{stl} & \texttt{true} & -- & -- & Write watertight STL solids \\
\texttt{stl\_facets} & 128 & -- & [8, 4096] & Azimuthal facets of the STL \\
\texttt{stl\_points}\,$^\ast$ & 220 & -- & [10, 100000] & Profile points of the candidate STLs \\
\texttt{stl\_points\_design}\,$^\ast$ & 400 & -- & [10, 100000] & Profile points of the design STL \\
\texttt{step} & \texttt{true} & -- & -- & Write a STEP solid of the design (needs the gmsh package) \\
\texttt{step\_points}\,$^\ast$ & 200 & -- & [10, 5000] & Profile points through which the STEP B-spline is fitted \\
\texttt{candidates} & \texttt{true} & -- & -- & Also export every profile of the sweep at the design fineness (STL) \\
\end{longtable}\addtocounter{table}{-1}}

\subsection*{\texttt{output}: Outputs}
\addcontentsline{toc}{subsection}{\texttt{output}: Outputs}
{\footnotesize
\begin{longtable}{@{}>{\raggedright\arraybackslash}p{3.3cm}>{\raggedright\arraybackslash}p{2.6cm}>{\raggedright\arraybackslash}p{1.3cm}>{\raggedright\arraybackslash}p{1.9cm}>{\raggedright\arraybackslash}p{6.1cm}@{}}
\toprule Parameter & Default & Unit & Valid & Meaning \\ \midrule \endhead \bottomrule \endfoot
\texttt{formats} & \texttt{[pdf, png]} & -- & pdf, png, svg & Figure file formats \\
\texttt{png\_dpi}\,$^\ast$ & 200 & dpi & [50, 1200] & Resolution of PNG figures \\
\texttt{plots} & \texttt{[all]} & -- & -- & Plots to draw: 'all' or a list of names (see the user guide) \\
\texttt{latex} & \texttt{false} & -- & -- & Also write LaTeX tables and macros (design-report style) \\
\texttt{text\_renderer} & \texttt{auto} & -- & auto, latex, mathtext & Plot text: 'latex' (needs a LaTeX installation), 'mathtext' (built in) or 'auto' \\
\texttt{plot\_options} & \{\,\} & -- & -- & Per-plot overrides, e.g. \{trajectory: \{xlim: [0, 40]\}\} \\
\end{longtable}\addtocounter{table}{-1}}



\subsection{Structured parameters}
\label{sec:structured}

\paragraph{\key{nose.profiles} and \key{nose.custom\_profiles}.} The sweep compares the profiles listed in
\key{nose.profiles}. The built-in names are those of the CFD database: \texttt{cone}, \texttt{ogive} (tangent),
\texttt{parabolic} ($K' = 1$), \texttt{power050} and \texttt{power075} (power series $n = \tfrac12$, $\tfrac34$),
\texttt{lvhaack} (Haack $C = \tfrac13$) and \texttt{vonkarman} (Haack $C = 0$). Other members of the same families
(and of any family added by a plugin) are defined in \key{custom\_profiles}:
\begin{lstlisting}[language=yaml]
nose:
  custom_profiles:
    power060: {family: power, n: 0.6, label: "Power $n=0.6$"}
    haack_c015: {family: haack, C: 0.15, label: "Haack $C=0.15$"}
  profiles: [vonkarman, power060, haack_c015]
\end{lstlisting}
The families and their parameters are listed by \texttt{python -m nosecone profiles}: \texttt{cone},
\texttt{tangent\_ogive}, \texttt{parabolic} ($K$), \texttt{power} ($n$), \texttt{haack} ($C$) and
\texttt{elliptical}. The \texttt{label} (LaTeX allowed) is used in plots and tables. A custom profile has no CFD
data until it is computed: set \key{aerodynamics.compute\_missing: true} (Section~\ref{sec:newcfd}).

\paragraph{\key{nose.fineness} and \key{nose.tip\_radius\_ratios}.} The sweep is the product
profiles $\times$ fineness $\times$ tip radii. The Weisshorn database covers $L/D = 3$--$5$ for sharp noses; a
fineness outside the computed range is accepted, but its wave drag is extrapolated with $C_{D,\rm w}(L/D)^2$ held
constant. Tip radii $r_n/R > 0$ need CFD data for the blunted nose at two fineness ratios or more (otherwise the
case is skipped, or computed if \key{compute\_missing} is set).

\paragraph{\key{heating.materials}.} One entry per material; add your own and refer to it by its key:
\begin{lstlisting}[language=yaml]
heating:
  materials:
    Ti64: {label: "Ti-6Al-4V", conductivity: 6.7, density: 4430, specific_heat: 526}
  tip_material: Ti64
\end{lstlisting}
\texttt{conductivity} in W/(m\,K), \texttt{density} in kg/m\textsuperscript{3}, \texttt{specific\_heat} in J/(kg\,K);
\texttt{label} is used in the plots and \texttt{short} (optional) in the LaTeX tables. Materials given in a file are
added to the default ones (CFRP, Al6082, PETCF).

\paragraph{\key{heating.shell\_stations}.} Positions along the shell where the wall temperature is computed:
\texttt{"0.05 m"} (distance from the tip) or \texttt{"0.5 L"} (fraction of the profile length).

\paragraph{\key{robustness.scenarios}.} Each scenario is a name plus any of the following changes to the nominal
model; the apogee of every profile is recomputed and shown relative to \key{robustness.reference\_profile}:
\begin{lstlisting}[language=yaml]
robustness:
  scenarios:
    - {name: nominal}
    - {name: "drag rise from M = 0.75", transonic_onset: 0.75}   # Mach of the start of the rise
    - {name: "transonic peak +20 %", transonic_factor: 1.2}      # wave drag x 1.2 below transonic_match
    - {name: "rest of vehicle drag +20 %", drag_factor_scale: 1.2}  # k x 1.2
    - {name: "LVO masses, 30 C day", mass_model: lvo, ground_temperature: 303.15}
\end{lstlisting}

\paragraph{\key{trade\_matrix}.} \key{weights} gives the relative weight of each criterion (\texttt{performance}:
apogee, \texttt{integration}: payload length inside the nose, \texttt{manufacture}, \texttt{stability}:
static-margin change). They are normalised to add up to one, so \texttt{\{performance: 4, integration: 2,
manufacture: 2, stability: 1\}} gives 0.44, 0.22, 0.22 and 0.11; a criterion left out of \key{weights} is left out
of the matrix. Performance, integration and stability are scored automatically from 1 (worst) to 5 (best) between
the worst and the best profile; \key{manufacture} is your judgement, one score per profile name. Only properties
of the hardware are scored: whether a profile has flown before is not a criterion.

\paragraph{\key{cfd.grids} and \key{cfd.studies}.} A grid level is \texttt{[cells along the nose, cells along the
cylinder, cells normal to the wall, first-cell height / R]}. A study is a list of CFD runs; each entry is solved for
every profile $\times$ fineness $\times$ tip radius, from the first Mach number down (each Mach number starts from
the previous solution):
\begin{lstlisting}[language=yaml]
cfd:
  studies:
    - {name: main, profiles: [power060], fineness: [3, 4, 5], machs: [1.6, 1.4, 1.2, 1.1]}
    - {name: blunt, profiles: [vonkarman], fineness: [4.5], tip_radius_ratios: [0.05, 0.1],
       grid: medium, machs: [1.6, 1.4, 1.2], save_fields: false}
\end{lstlisting}
Optional keys of a study: \texttt{tip\_radius\_ratios} (default \texttt{[0]}), \texttt{grid} (default
\texttt{medium}), \texttt{save\_fields} (save the flow fields at \key{cfd.save\_fields} Mach numbers) and
\texttt{max\_iter} (iteration limit of this study, instead of \key{cfd.max\_iter}). Results go to
\key{cfd.\allowbreak cp\_dir}; \texttt{python -m nosecone rebuild-database} turns them into the
database CSV. The study name becomes the \texttt{study} column of the database; \key{aerodynamics.\allowbreak studies}
chooses which studies feed
the wave-drag model (\texttt{main}: the sweep; \texttt{f45}: direct runs at $L/D = 4.5$).

\paragraph{\key{output.plots} and \key{output.plot\_options}.} \key{plots} is \texttt{[all]} or a list of plot
names: \texttt{profiles}, \texttt{schematics}, \texttt{validation}, \texttt{cdw}, \texttt{cp}, \texttt{trajectory},
\texttt{breakdown}, \texttt{trade}, \texttt{robustness}, \texttt{bluntness}, \texttt{thermal}, \texttt{final},
\texttt{gallery}, \texttt{fields}, \texttt{grid}, \texttt{renders} (Section~\ref{sec:plots}). Plots that need data
the study does not have (e.g.\ \texttt{validation} without cone runs) are skipped with a message. \key{plot\_options}
overrides cosmetic settings per plot; the axis ranges that depend on the vehicle are automatic unless set:
\begin{lstlisting}[language=yaml]
output:
  plot_options:
    trajectory: {xlim: [0, 40], ylim: [0, 2.1]}     # time [s] and Mach axes
    thermal: {xlim: [0, 30], ylim: [0, 175]}        # time [s] and temperature [degC] axes
\end{lstlisting}

\begin{table}[H]
\centering
\caption{Options of \key{output.plot\_options} (defaults in \texttt{PLOT\_DEFAULTS}, \texttt{nosecone/plotting/figures.py}).
``auto'': chosen from the data. Lengths in mm scale with the diameter.}
\label{tab:plotopts}
\footnotesize
\begin{tabular}{@{}>{\raggedright\arraybackslash}p{2.2cm}>{\raggedright\arraybackslash}p{13.8cm}@{}}
\toprule
Plot & Options (default) \\
\midrule
\texttt{profiles} & \texttt{fineness} (baseline), \texttt{ylim\_mm} (100), \texttt{tip\_xlim\_mm} (60), \texttt{tip\_ylim\_mm} (38), label positions \\
\texttt{cdw} & \texttt{fineness} (baseline), \texttt{mach\_range} (0.6, 1.65), \texttt{ylim} (0, 0.125), \texttt{fineness\_ylim} (0.03, 0.14), \texttt{fineness\_mach} (1.4) \\
\texttt{cp} & \texttt{fineness} (baseline), \texttt{machs} (1.1, 1.4), \texttt{xlim} (0, 1.35), \texttt{ylim} ($-0.25$, 0.45) \\
\texttt{trajectory} & \texttt{xlim}, \texttt{ylim} (auto), \texttt{mach\_bins\_max} (auto), \texttt{bin\_width} (0.1) \\
\texttt{breakdown} & \texttt{altitude} (3000 m), \texttt{mach\_range} (0.1 to the highest tabulated Mach), \texttt{annotate\_mach} (1.4) \\
\texttt{trade} & \texttt{payload\_ylim} (0 to 1.5 payload lengths) \\
\texttt{bluntness} & \texttt{machs} (1.2, 1.4, 1.6), \texttt{guide} ([kW/(m\textsuperscript{2}K), m] of the $r_n^{-1/2}$ guide line; auto) \\
\texttt{thermal} & \texttt{xlim}, \texttt{ylim} (auto) \\
\texttt{fields} & \texttt{machs} (1.4), \texttt{xlim} (auto), \texttt{ylim} ($\pm3.2\,R$), \texttt{clim} ($-0.25$, 0, 0.45) \\
\texttt{grid}, \texttt{schematics} & \texttt{level} (coarse); \texttt{fineness} (2, drawing only) \\
\texttt{validation}, \texttt{robustness}, \texttt{gallery} & Mach range, colours and markers; \texttt{gallery}: \texttt{mach} (1.4) \\
\bottomrule
\end{tabular}
\end{table}
